% docs/.paper/sec_related.tex
\section{Related Work}
\label{sec:related}

\subsection{Rule-Based Transpilation \& AST Rewriting}
\label{sec:related_transpilers}

Traditional code translation systems rely on deterministic abstract syntax tree (AST) transformations. \textsc{C2Rust}~\cite{c2rust,p38_c2rust} translates C code into Rust by preserving C memory layouts, emitting \texttt{unsafe} blocks and raw pointer operations. While semantics are preserved on supported syntax subsets, the emitted code requires extensive human refactoring to achieve idiomatic safety and maintainability~\cite{siddiq2024oxidizer,p05_oxidizer,p50_rav1d_c2rust_case,p51_libexpat_c2rust_case}. Rule-based tools like \textsc{CXGO}, \textsc{OpenRewrite}~\cite{p39_openrewrite}, and \textsc{JavaC2Rust}~\cite{p59_javac2rust_tooling} use specialized pattern rewrites. However, they struggle to generalize across divergent type systems, complex macro expansions~\cite{p57_proc_macro_transpilation}, and foreign function interfaces~\cite{p56_rust_ffi_verification}.

\subsection{Neural Code Translation \& Program Synthesis}
\label{sec:related_neural}

Neural approaches formulate source code translation as sequence-to-sequence language modeling. \textsc{TransCoder}~\cite{roziere2020transcoder} introduced unsupervised machine translation for programming languages using cross-lingual masked modeling. Follow-up works like \textsc{AlphaTrans}~\cite{yang2024alphatrans,p02_alphatrans} introduced modular call-graph analysis and iterative validation feedback. Dedicated code models such as \textsc{DeepSeek-Coder}~\cite{p18_deepseekcoder,p44_deepseek_coder_v2}, \textsc{StarCoder2}~\cite{p22_starcoder2}, \textsc{CodeLlama}~\cite{p32_codellama}, \textsc{Phi-3}~\cite{p27_phi3,p45_phi_evolution}, and \textsc{Qwen2.5-Coder}~\cite{p21_qwen2_5coder,p41_qwen_moe} have achieved high accuracy on isolated benchmark functions. However, these models operate primarily on single files, lacking whole-repository dependency tracking and cross-file context orchestration.

\subsection{Multi-Agent Software Engineering Systems}
\label{sec:related_mas}

Multi-agent software engineering frameworks coordinate specialized language model agents across discrete development roles. \textsc{MetaGPT}~\cite{hong2023metagpt,p11_metagpt,p42_metagpt_deepread} incorporates standard operating procedures to structure role-based artifact handoffs. \textsc{ChatDev}~\cite{qian2023chatdev,p12_chatdev,p46_chatdev_deepread} introduces communicative de-hallucination loops for iterative code refinement. \textsc{HyperAgent}~\cite{p13_hyperagent,p47_hyperagent_deepread} and \textsc{AgentVerse}~\cite{p26_agentverse} demonstrate multi-agent collaboration for scalable repository navigation and dynamic recruitment. Frameworks such as \textsc{AutoGen}~\cite{p16_autogen,p48_autogen_deepread}, \textsc{CAID}~\cite{p54_caid}, and \textsc{OpenHands}~\cite{p53_openhands_deepread} explore asynchronous blackboards and human-in-the-loop checkpointing. 

In repository-scale translation, \textsc{ReCodeAgent}~\cite{ibrahimzada2026recodeagent,p01_recodeagent} demonstrated that multi-agent decomposition significantly improves compilation rates over monolithic prompting. Retrieval-augmented architectures like \textsc{CodePlan}~\cite{p03_codeplan}, \textsc{RepoCoder}~\cite{p61_repocoder}, and \textsc{RepoGraph}~\cite{p62_repograph} leverage structured dependency graphs to supply precise context slices. MAgHARCM builds upon these foundations to prove that repository-scale translation can execute entirely on local, quantized Small Language Models through reverse-topological dependency scheduling and bounded context navigation.
